\subsection{অবৈধ ইনপুট, validity ও decimal-to-BCD}
\chThreeExample{১}{একাধিক ইনপুটে ভুল পরিচয়}
সাধারণ 4-to-2 এনকোডারে $D_1=D_2=1$।
\chThreeAnswer $A=D_2+D_3=1$, $B=D_1+D_3=1$; আউটপুট 11, যা $D_3$ সক্রিয় থাকার কোড। কিন্তু $D_3$ সক্রিয় নয়। একটিমাত্র ইনপুট সক্রিয় থাকার শর্ত বাদ দিলে ভুল পরিচয় পাওয়া যায়।

\chThreeExample{২}{কোনো ইনপুট নেই, নাকি $D_0$?}
সকল ইনপুট ০ হলে আউটপুট 00; $D_0=1$ হলেও 00।
\chThreeAnswer আলাদা বৈধতা সংকেত $V=D_0+D_1+D_2+D_3$ রাখা যায়। $V=0$ হলে কোনো ইনপুট সক্রিয় নয়; $V=1,AB=00$ হলে একটিমাত্র ইনপুটের শর্তে $D_0$ সক্রিয়।

\chThreeExample{৩}{অগ্রাধিকার নির্বাচন}
8-to-3 priority encoder-এ বড় নম্বরের ইনপুটের অগ্রাধিকার বেশি। $D_2,D_5,D_7$ একসঙ্গে সক্রিয়।
\chThreeAnswer সর্বোচ্চ অগ্রাধিকার $D_7$; আদর্শ active-high কোড 111। বাস্তব IC-তে active-low ইনপুট/আউটপুটও থাকতে পারে; তখন datasheet অনুযায়ী polarity নিতে হয়।

\chThreeExample{৪}{দশমিক অঙ্ক থেকে BCD}
একটি decimal-to-BCD encoder-এ $D_0$ থেকে $D_9$ পর্যন্ত দশটি ইনপুট। $D_9$ একা সক্রিয় হলে আউটপুট $Q_3Q_2Q_1Q_0=1001$।
\chThreeAnswer
\[\begin{aligned}Q_3&=D_8+D_9,\\Q_2&=D_4+D_5+D_6+D_7,\\Q_1&=D_2+D_3+D_6+D_7,\\Q_0&=D_1+D_3+D_5+D_7+D_9.\end{aligned}\]
এখানে চার আউটপুট থাকলেও ১৬ ইনপুট নেই। সাধারণ $2^n$-to-$n$ binary encoder-এর নিয়ম এবং নির্দিষ্ট decimal encoder-এর ব্যবহৃত ইনপুট পৃথক করে বুঝতে হবে।

% Mahbub teaching-type follow-up, 2026-10-14
\chThreeExample{৫}{output দেখে সক্রিয় input ও রেখার সংখ্যা}
একটি সাধারণ binary encoder-এর output তিন বিট এবং কোড 101। input-এ একটিমাত্র রেখা active-high। কয়টি input এবং কোনটি সক্রিয়?
\chThreeAnswer $2^3=8$টি input: $D_0$ থেকে $D_7$। $101_2=5$, তাই $D_5=1$। একই কোডের MSB/LSB উল্টো নাম দিলে পরিচয় বদলাতে পারে; output order প্রশ্নে নির্দিষ্ট করতে হয়। একটিমাত্র active input-এর শর্ত না থাকলে সাধারণ encoder থেকে এই সিদ্ধান্ত নিশ্চিত নয়।
